\documentclass[11pt,a4paper]{article}
\usepackage[utf8]{inputenc}
\usepackage[T1]{fontenc}
\usepackage{geometry}
\usepackage{graphicx}
\usepackage{hyperref}
\usepackage{listings}
\usepackage{xcolor}
\usepackage{amsmath}

\geometry{margin=2.5cm}

\title{AI-Driven Distributed Fog Load Balancing \& Anomaly-Aware Routing\\
Technical Specification}
\author{Engineering Project}
\date{\today}

\begin{document}

\maketitle

\section{Project Overview}

\subsection{Title}
AI-Driven Distributed Fog Load Balancing \& Anomaly-Aware Routing

\subsection{Objective}
Implement a real distributed fog computing architecture demonstrating mastery of fog computing, networking protocols, and artificial intelligence through concrete implementation, measurable performance metrics, and architectural comparisons.

\subsection{Key Requirements}
\begin{itemize}
    \item Real distributed fog computing architecture (not conceptual simulation)
    \item Multiple fog nodes in distributed mesh topology
    \item SDN-based dynamic routing (OpenFlow + Ryu)
    \item AI components embedded in each fog node:
    \begin{itemize}
        \item LSTM for future load prediction
        \item Autoencoder for network anomaly detection
        \item Reinforcement Learning for intelligent routing decisions
    \end{itemize}
    \item Comparison of three architectures:
    \begin{enumerate}
        \item Centralized fog (single controller, static routing, no AI)
        \item Distributed fog without AI (rule-based load balancing)
        \item Distributed fog with AI (adaptive, predictive, anomaly-aware)
    \end{enumerate}
\end{itemize}

\section{System Architecture}

\subsection{Network Topology}
\begin{itemize}
    \item \textbf{Topology Type}: Distributed mesh network
    \item \textbf{Fog Nodes}: Minimum 4-6 nodes for meaningful distributed behavior
    \item \textbf{IoT Devices}: Multiple IoT devices generating traffic
    \item \textbf{Cloud Layer}: Used only for monitoring and historical analysis (excluded from real-time decision-making)
    \item \textbf{Inter-node Communication}: gRPC for fog-to-fog communication
    \item \textbf{IoT Communication}: MQTT protocol
\end{itemize}

\subsection{Technology Stack}

\subsubsection{Network Emulation}
\begin{itemize}
    \item \textbf{Fogbed}: Container-based network emulation (Containernet + Mininet)
    \item \textbf{Open vSwitch}: Virtual switch for SDN
    \item \textbf{OpenFlow}: Protocol version 1.3 or 1.5
    \item \textbf{Ryu SDN Controller}: Python-based SDN controller
\end{itemize}

\subsubsection{Communication Protocols}
\begin{itemize}
    \item \textbf{MQTT}: IoT device messaging (Eclipse Mosquitto broker)
    \item \textbf{gRPC}: Inter-fog node communication (Protocol Buffers)
    \item \textbf{REST API}: Optional for management interfaces
\end{itemize}

\subsubsection{AI/ML Components}
\begin{itemize}
    \item \textbf{LSTM}: PyTorch or TensorFlow for load prediction
    \item \textbf{Autoencoder}: PyTorch or TensorFlow for anomaly detection
    \item \textbf{Reinforcement Learning}: Stable-Baselines3 or custom DQN implementation
\end{itemize}

\subsubsection{Monitoring \& Metrics}
\begin{itemize}
    \item \textbf{Prometheus}: Time-series metrics collection
    \item \textbf{Grafana}: Visualization and dashboards
    \item \textbf{Custom Exporters}: For fog node metrics (CPU, memory, network, latency)
\end{itemize}

\subsubsection{Programming Languages}
\begin{itemize}
    \item \textbf{Python 3.8+}: Primary language for controllers, AI components, and orchestration
    \item \textbf{Bash}: Scripts for deployment and automation
\end{itemize}

\section{Component Specifications}

\subsection{Fog Node Components}

Each fog node implements:

\subsubsection{Load Monitor}
\begin{itemize}
    \item Collects real-time metrics: CPU usage, memory usage, network bandwidth, active connections
    \item Sampling interval: 1-5 seconds
    \item Exports metrics to Prometheus
\end{itemize}

\subsubsection{LSTM Load Predictor}
\begin{itemize}
    \item Input: Historical load metrics (time window: 60-300 seconds)
    \item Output: Predicted load for next 10-30 seconds
    \item Architecture: 2-3 LSTM layers, dropout for regularization
    \item Training: Online learning with periodic retraining
    \item Update frequency: Every 10-30 seconds
\end{itemize}

\subsubsection{Autoencoder Anomaly Detector}
\begin{itemize}
    \item Input: Network traffic features (packet rate, flow count, latency, error rate)
    \item Architecture: Encoder-decoder with bottleneck layer
    \item Threshold: Reconstruction error threshold for anomaly detection
    \item Output: Anomaly score (0-1) and binary classification
    \item Update frequency: Real-time (per packet batch)
\end{itemize}

\subsubsection{RL Routing Agent}
\begin{itemize}
    \item State space: Current node load, predicted load, neighbor node states, anomaly scores
    \item Action space: Route selection (next hop selection)
    \item Reward function: Combines latency, load balance, and anomaly avoidance
    \item Algorithm: DQN (Deep Q-Network) or PPO (Proximal Policy Optimization)
    \item Update frequency: Per routing decision (every packet or flow)
\end{itemize}

\subsubsection{gRPC Service}
\begin{itemize}
    \item Service definition: Node state exchange, routing queries, load information
    \item Protocol: Protocol Buffers (protobuf)
    \item Communication: Asynchronous or synchronous based on requirements
\end{itemize}

\subsection{SDN Controller (Ryu)}

\subsubsection{OpenFlow Applications}
\begin{itemize}
    \item \textbf{Flow Manager}: Installs/modifies flow rules based on routing decisions
    \item \textbf{Topology Discovery}: Discovers network topology using LLDP
    \item \textbf{Statistics Collector}: Collects flow statistics for monitoring
    \item \textbf{Routing Application}: Integrates with fog node AI for routing decisions
\end{itemize}

\subsubsection{Controller Modes}
\begin{enumerate}
    \item \textbf{Centralized Mode}: Single controller makes all routing decisions (Architecture 1)
    \item \textbf{Distributed Mode}: Each fog node makes local routing decisions (Architectures 2 \& 3)
\end{enumerate}

\subsection{IoT Traffic Generator}

\subsubsection{Traffic Patterns}
\begin{itemize}
    \item \textbf{Periodic}: Regular interval traffic (sensors)
    \item \textbf{Burst}: Sudden traffic spikes (events)
    \item \textbf{Variable}: Random intervals with different intensities
    \item \textbf{Anomalous}: Injected anomalies for testing (DDoS-like, unusual patterns)
\end{itemize}

\subsubsection{MQTT Topics}
\begin{itemize}
    \item \texttt{iot/devices/\{device\_id\}/data}: Sensor data
    \item \texttt{iot/devices/\{device\_id\}/status}: Device status
    \item \texttt{fog/nodes/\{node\_id\}/routing}: Routing updates
\end{itemize}

\section{Architecture Variants}

\subsection{Architecture 1: Centralized Fog}
\begin{itemize}
    \item \textbf{Controller}: Single Ryu controller
    \item \textbf{Routing}: Static or simple shortest-path routing
    \item \textbf{AI}: None
    \item \textbf{Decision Making}: Centralized in controller
    \item \textbf{Load Balancing}: Round-robin or static assignment
\end{itemize}

\subsection{Architecture 2: Distributed Fog (No AI)}
\begin{itemize}
    \item \textbf{Controllers}: Multiple Ryu controllers (one per fog node or shared)
    \item \textbf{Routing}: Rule-based load balancing
    \item \textbf{AI}: None
    \item \textbf{Decision Making}: Distributed, based on local rules
    \item \textbf{Load Balancing}: Weighted round-robin, least-loaded, or similar heuristics
\end{itemize}

\subsection{Architecture 3: Distributed Fog (With AI)}
\begin{itemize}
    \item \textbf{Controllers}: Multiple Ryu controllers with AI integration
    \item \textbf{Routing}: AI-driven adaptive routing
    \item \textbf{AI Components}: LSTM, Autoencoder, RL agent active
    \item \textbf{Decision Making}: Distributed, AI-enhanced
    \item \textbf{Load Balancing}: Predictive and anomaly-aware
\end{itemize}

\section{Performance Metrics}

\subsection{Primary Metrics}
\begin{itemize}
    \item \textbf{Latency}: End-to-end packet latency (mean, p50, p95, p99)
    \item \textbf{Throughput}: Packets/second, bits/second
    \item \textbf{Load Balance}: Standard deviation of node loads
    \item \textbf{Anomaly Detection}: Precision, recall, F1-score, detection time
    \item \textbf{Resource Utilization}: CPU, memory, network bandwidth per node
\end{itemize}

\subsection{Secondary Metrics}
\begin{itemize}
    \item \textbf{Packet Loss}: Percentage of dropped packets
    \item \textbf{Flow Completion Time}: Time for complete flows
    \item \textbf{Convergence Time}: Time for system to adapt to changes
    \item \textbf{AI Model Accuracy}: LSTM prediction error, Autoencoder reconstruction error
    \item \textbf{RL Performance}: Reward over time, policy convergence
\end{itemize}

\section{Test Scenarios}

\subsection{Scenario 1: Baseline Performance}
\begin{itemize}
    \item Steady-state traffic
    \item All three architectures
    \item Measure baseline latency, throughput, load balance
\end{itemize}

\subsection{Scenario 2: Traffic Burst}
\begin{itemize}
    \item Sudden traffic increase (2x, 5x, 10x)
    \item Measure adaptation time and performance degradation
    \item Compare predictive capabilities of Architecture 3
\end{itemize}

\subsection{Scenario 3: Node Failure}
\begin{itemize}
    \item Random node failure
    \item Measure recovery time and routing adaptation
    \item Test fault tolerance
\end{itemize}

\subsection{Scenario 4: Anomaly Injection}
\begin{itemize}
    \item Inject network anomalies (DDoS, unusual patterns)
    \item Measure detection accuracy and time
    \item Test routing adaptation to avoid anomalies
\end{itemize}

\subsection{Scenario 5: Scalability}
\begin{itemize}
    \item Vary number of fog nodes (4, 6, 8, 10)
    \item Vary number of IoT devices (10, 50, 100, 200)
    \item Measure performance scaling
\end{itemize}

\subsection{Scenario 6: Dynamic Load}
\begin{itemize}
    \item Variable traffic patterns over time
    \item Measure long-term performance and stability
    \item Test AI model adaptation
\end{itemize}

\section{Implementation Structure}

\subsection{Directory Structure}
\begin{verbatim}
project/
├── technical_sheet.tex
├── README.md
├── requirements.txt
├── docker-compose.yml
├── scripts/
│   ├── setup.sh
│   ├── deploy.sh
│   └── cleanup.sh
├── fogbed/
│   ├── topology.py
│   └── network_config.py
├── sdn_controller/
│   ├── ryu_apps/
│   │   ├── centralized_controller.py
│   │   ├── distributed_controller.py
│   │   └── ai_routing_app.py
│   └── ofproto/
│       └── fog_proto.py
├── fog_nodes/
│   ├── node_base.py
│   ├── load_monitor.py
│   ├── lstm_predictor.py
│   ├── anomaly_detector.py
│   ├── rl_agent.py
│   └── grpc_service.py
├── iot/
│   ├── traffic_generator.py
│   └── mqtt_client.py
├── monitoring/
│   ├── prometheus/
│   │   └── prometheus.yml
│   ├── grafana/
│   │   └── dashboards/
│   └── exporters/
│       └── fog_exporter.py
├── ai_models/
│   ├── lstm/
│   │   ├── model.py
│   │   └── train.py
│   ├── autoencoder/
│   │   ├── model.py
│   │   └── train.py
│   └── rl/
│       ├── dqn_agent.py
│       └── environment.py
├── tests/
│   ├── test_scenarios.py
│   └── performance_tests.py
└── results/
    ├── metrics/
    └── logs/
\end{verbatim}

\section{Deployment Requirements}

\subsection{System Requirements}
\begin{itemize}
    \item \textbf{OS}: Linux (Ubuntu 20.04+ recommended)
    \item \textbf{Python}: 3.8 or higher
    \item \textbf{Docker}: For containerization (optional but recommended)
    \item \textbf{Resources}: Minimum 8GB RAM, 4 CPU cores (more for larger topologies)
\end{itemize}

\subsection{Dependencies}
\begin{itemize}
    \item Fogbed/Containernet
    \item Ryu SDN Framework
    \item PyTorch or TensorFlow
    \item Stable-Baselines3 (for RL)
    \item Prometheus, Grafana
    \item MQTT broker (Mosquitto)
    \item gRPC and Protocol Buffers
    \item NumPy, Pandas, Scikit-learn
\end{itemize}

\section{Evaluation Criteria}

\subsection{Technical Correctness}
\begin{itemize}
    \item Proper SDN implementation with OpenFlow
    \item Correct fog node communication (gRPC)
    \item Functional AI components
    \item Accurate metrics collection
\end{itemize}

\subsection{Performance Comparison}
\begin{itemize}
    \item Quantitative comparison of three architectures
    \item Statistical significance of results
    \item Clear performance trade-offs identified
\end{itemize}

\subsection{Code Quality}
\begin{itemize}
    \item Well-structured, modular code
    \item Documentation and comments
    \item Reproducible experiments
    \item Error handling and logging
\end{itemize}

\section{Timeline \& Milestones}

\begin{enumerate}
    \item \textbf{Week 1-2}: Environment setup, Fogbed topology implementation
    \item \textbf{Week 3-4}: SDN controller implementation (all three variants)
    \item \textbf{Week 5-6}: AI components implementation (LSTM, Autoencoder, RL)
    \item \textbf{Week 7-8}: Integration and testing
    \item \textbf{Week 9-10}: Performance evaluation and comparison
    \item \textbf{Week 11-12}: Documentation and finalization
\end{enumerate}

\section{Deliverables}

\begin{itemize}
    \item Complete source code with documentation
    \item Technical specification (this document)
    \item Performance evaluation report with metrics
    \item Deployment instructions and scripts
    \item Test scenarios and results
    \item Comparison analysis of three architectures
\end{itemize}

\end{document}

